\documentclass[10pt,a4paper]{amsart}
\usepackage[margin=19mm]{geometry}
\usepackage{amsmath,amssymb,mathtools}
\usepackage[scr=rsfs,cal=euler]{mathalfa}
\usepackage{xfrac}
\usepackage[unicode,hidelinks,bookmarks=false]{hyperref}
\newcommand{\ie}{i.e.\ }
\newcommand{\cf}{cf.\ }
\newcommand{\resp}{respectively\ }
\newcommand{\Subsection}{\subsection}
\providecommand{\nobreakdash}{\nobreak\hbox{-}}
\setlength{\emergencystretch}{2em}
\multlinegap0pt
\usepackage{siunitx}
\usepackage{siunitx}
\usepackage{siunitx}
\usepackage{siunitx}
\usepackage{amssymb}
\usepackage{mathtools}
\usepackage{siunitx}
\usepackage{xcolor}
\usepackage{float}
\usepackage{algorithm}\usepackage{algpseudocode}
\usepackage[shortlabels]{enumitem}
\newtheorem{theorem}{Theorem}
\title{A Mathematical Model to Predict Growth and Treatment for UPS Cancer}
\author{Sumit Roy}
\date{Source: arXiv:2512.21620v3 (CC BY 4.0)}
\begin{document}
\maketitle

\noindent\emph{Source and licence.} A Mathematical Model to Predict Growth and Treatment for UPS Cancer. Version arXiv:2512.21620v3 (CC BY 4.0), distributed by arXiv under the Creative Commons Attribution 4.0 International licence, http://creativecommons.org/licenses/by/4.0/, as recorded in the arXiv licence metadata for this exact version and shown on the abstract page https://arxiv.org/abs/2512.21620v3. Full licence notice: \texttt{licenses/arxiv-2512.21620-CC-BY-4.0.txt}.\par\smallskip

\noindent\textit{Editorial scope.} Counted unit: the article's theorem thm:rt (source0.tex lines 894--919). The article's complete original proof of it is delivered with it (source0.tex lines 921--981, one \texttt{proof} environment of 61 lines). No mathematics is written, repaired or re-derived: the statement and the proof are the article's own lines, in the article's own notation, and no retained line is substituted or re-flowed.  No further source-local context is needed: every cross-reference of every retained line resolves inside the two delivered spans.  Nothing was omitted from any retained span, so the package carries no omission record.  No retained line contains a citation, so the wrapper carries no bibliography and no bibliography entry.  No retained line contains a figure environment, an image-inclusion command or drawing code, so no image is delivered and the package declares no asset dependency.  Editorial formatting only, in addition: an AMS \texttt{amsart} preamble that restates the article's own declarations and macros used by the retained lines, namely the environments \texttt{theorem} on separate counters, together with the standard packages those lines need.  The retained environments print in this document as Theorem 1 in order of appearance, whereas the article prints the same items with the numbering of its own full document.  The complete native source of the article as distributed in the arXiv e-print for this exact version is bundled unchanged as \texttt{sources/191674/source0.tex}.
 \begin{theorem}\label{thm:rt}
Consider the radiation-induced volume decay
\[
\frac{dV}{dt} = -\Psi(t)D(t)V,
\]
where
\[
\Psi(t) = \alpha + \beta \sin\!\left(\frac{2\pi t}{\tau}\right)
\]
denotes the time-dependent radiosensitivity and assume $\Psi(t)\ge0$ for $t\in[0,T]$. Let the admissible dose schedules satisfy
\[
0 \le D(t) \le D_{\max}, 
\qquad 
\int_0^T D(t)\,dt \le D_{\mathrm{total}}.
\]
Then any dose schedule $D^*(t)$ minimizing the terminal tumor volume $V(T)$ is of bang--bang type. More precisely, there exists a constant threshold $\lambda$ such that
\[
D^*(t) =
\begin{cases}
D_{\max}, & \text{if } \Psi(t) > \lambda,\\[2mm]
0, & \text{if } \Psi(t) < \lambda,
\end{cases}
\]
for almost every $t\in[0,T]$. On the set where $\Psi(t)=\lambda$, the control may take intermediate values so that the total dose
constraint is satisfied.
\end{theorem}

\begin{proof}
For any admissible control $D(t)$, the state equation can be solved explicitly. Integrating the equation gives
\[
V(T) = V(0)\exp\!\left(-\int_0^T \Psi(t)D(t)\,dt\right).
\]
Since $V(0)$ is fixed and the exponential function is strictly decreasing in its exponent, minimizing the terminal tumor volume
$V(T)$ is equivalent to maximizing the functional
\[
J(D) = \int_0^T \Psi(t)D(t)\,dt
\]
over all admissible controls.

Suppose that $D^*(t)$ is an optimal control. We show that the radiation dose must be concentrated at times where the radiosensitivity $\Psi(t)$ is largest. Assume that there exist measurable sets $A,B\subset[0,T]$ of positive measure such that
\[
\Psi(t)>\Psi(s) \quad \text{for all } t\in A,\ s\in B,
\]
and
\[
D^*(t) < D_{\max} \quad \text{for } t\in A, 
\qquad
D^*(s) > 0 \quad \text{for } s\in B.
\]

Choose $\varepsilon>0$ sufficiently small so that the modified control
\[
\widetilde D(t) =
\begin{cases}
D^*(t)+\varepsilon, & t\in A,\\
D^*(t)-\varepsilon, & t\in B,\\
D^*(t), & \text{otherwise},
\end{cases}
\]
remains admissible. By choosing the sets $A$ and $B$ with equal measure, the total dose constraint is preserved.

The corresponding change in the objective functional is
\[
J(\widetilde D) - J(D^*)
= \varepsilon
\left(
\int_A \Psi(t)\,dt
-
\int_B \Psi(t)\,dt
\right).
\]
Since $\Psi(t)>\Psi(s)$ for $t\in A$ and $s\in B$, the right-hand side is positive, contradicting the optimality of $D^*$.

Therefore such sets cannot exist. Consequently, if radiation is administered at some time $t$, then it must be delivered at the
maximum allowable rate during periods where $\Psi(t)$ is larger. This implies the existence of a threshold $\lambda$ such that
\[
D^*(t)=D_{\max} \quad \text{when } \Psi(t)>\lambda,
\qquad
D^*(t)=0 \quad \text{when } \Psi(t)<\lambda.
\]

Finally, the value of $\lambda$ is determined by the total dose constraint. If the entire dose budget is used, $\lambda$ is chosen so that
\[
\int_0^T D^*(t)\,dt = D_{\mathrm{total}}.
\]
If the dose budget satisfies $D_{\mathrm{total}} \ge T D_{\max}$, the constraint is inactive and the optimal schedule is simply
$D^*(t) = D_{\max}$ for all $t\in[0,T]$.
\end{proof}

\end{document}
